Пусть $E_1, E_2$ --- линейные нормированные пространства над полем $\mathbb{K}$ ($\R$ или $\C$).

$A: E_1 \to E_2$ будем называть \textit{оператором}, а 
$f: E \to \mathbb{K}$ --- функционалом.

\begin{definition}
Оператор $A$ называется \textit{ограниченным}, если для любого ограниченного $M \subset E_1$
образ $A(M)$ ограничен в $E_2$.
\end{definition}

\begin{note}
Ограниченность оператора не означает, что множество его значений ограничено.
Например, тождественный оператор $I: E \to E$ является ограниченным, но множество его значений (всё $E$) неограничено.
\end{note}

\begin{note}
Есть альтернативные определения ограниченности, которые эквивалентны для линейных операторов и доказаны в утверждении \ref{equiv-bounded-oper}. Они будут так же использоваться далее и в этом билете.
\end{note}

\begin{definition}
Для линейных операторов можно ввести следующие определения:

Образ оператора: $\im A = \{ y \in E_2 : \exists x \in E_1 \ Ax = y\}$

Ядро оператора: $\ke A = \{ x \in E_1 : Ax = 0\}$
\end{definition}

\begin{note}
$\im A$ --- линейное многообразие, но не обязательно подпространство.
$\ke A$ --- линейное многообразие.
\end{note}

\begin{definition}
Линейный оператор $A$ называется \textit{непрерывным}, 
если для любой последовательности $x_n \to x$ выполнено $Ax_n \to Ax$.
\end{definition}

\begin{theorem}
Пусть $E_1, E_2$ --- линейные нормированные пространства. $A: E_1 \to E_2$ --- линейный оператор.
Тогда $A$ --- ограниченный тогда и только тогда, когда $A$ --- непрерывный.
\end{theorem}

\begin{proof}\\
\textit{\imp{ограниченность}{непрерывность}}\\
Равносильно: $x_n \to x \Leftrightarrow \|x_n - x\| \to 0$.
\[
\|Ax_n - Ax\|_{E_2} = \|A(x_n - x)\|_{E_2} \leqslant \|A\| \cdot \|x_n - x\|_{E_1} \to 0.
\]

\noindent\textit{\imp{непрерывность}{ограниченность}}\\
Предположим противное, то есть $A$ не является ограниченным. 
$\forall K\ \exists x : \|Ax\|_{E_2} > K \|x\|_{E_1}$. Пусть $K$ пробегает все натуральные числа, тогда образуется последовательность $\{x_n\}$ такая, что $\|Ax_n\|_{E_2} > n \|x_n\|_{E_1}$. Все $x_n$, очевидно, ненулевые.
Рассмотрим последовательность $y_n = \frac{1}{n} \frac{x_n}{\|x_n\|_{E_1}} \to 0$.

\[
\|A y_n\|_{E_2} = \frac{\|A x_n\|_{E_2}}{n \cdot \|x_n\|_{E_1}} > 1\ \forall n
\]

Но из-за непрерывности оператора $A y_n \to A 0 = 0$. Противоречие.
\end{proof}